\section{Experimental Setup}
\label{sec:setup}

\noindent\textbf{Sensors and data.} Sensor A is a GelSlim-type fingertip with a $17.5\,$mm gel view and six colored LED emitters. A robot presses it against 12 ridge-pattern objects with known geometry. Ground-truth (GT) depth and normals are mesh projections at the recorded pose, rather than measurements of membrane deformation. After filtering, $4{,}669$ frames remain; every tenth frame is assigned to the validation split ($472$ frames). A second session, recorded on another day after the objects were re-placed, contains $854$ reviewed test frames from 11 objects. Their poses are re-estimated from contact masks. Sensor B, \sensorB{}~\cite{sipos2024gelslim4}, has a marker array, different illumination and a $12\,$mm view. We use its public dataset of the same 12 objects ($1{,}620$ test frames). We re-register its clipped, smoothed depth labels to the meshes by pose fitting ($0.018\,$mm mean residual) and use the resulting targets for every sensor-B method.

\noindent\textbf{Label budget.} Each object contributes $\Kreal\in\{10,25,50,100\}$ evenly spaced labeled training frames. The teacher, renderer and aligner share these frames, which represent $3\%$, $7\%$, $14\%$ and $29\%$ of the sensor-A training set. All other training images are available without labels. Our primary setting is $K\le25$; the larger budgets show how performance changes as more labels become available. Validation and test images are excluded from component training. Validation labels are used for checkpoint selection and are not included in the stated training-label budget.

\noindent\textbf{Models and training.} All methods use a frozen SITR-B/18 encoder~\cite{gupta2025sitr}, a DPT decoder~\cite{ranftl2021dpt} for depth and normals, and a planar-pose MLP ($12$M trainable parameters). We train the main sensor-A models for 150 epochs and use 50 epochs for ablations, three-seed checks and sensor B. Aligner perturbations span $\pm8^\circ$, $\pm1.5\,$mm in-plane and $\pm0.15\,$mm in indentation; half use ranges scaled by $0.3$. Training uses smooth L1 for 8,000 steps (batch 64, AdamW, learning rate $3\times10^{-4}$). Main sensor-A batches mix 24 real and 40 synthetic examples; renderer-free and sensor-B aligners use real examples only. The main sensor-A teacher uses the $K$ real frames and learned-renderer data, whereas the other two use a real-only teacher. The sensor-A student additionally uses synthetic images from a renderer retrained on labeled and aligned pairs.

\noindent\textbf{Renderers and cost.} The physical renderer uses Blender/Cycles with illumination fitted to a no-contact image and manually tuned contact parameters. The learned renderer maps depth to RGB with a pix2pix network~\cite{isola2017pix2pix}. Both use the same $5{,}345$ simulated contacts. An additional \emph{leaky renderer} reference uses all $4{,}197$ real training pairs: it exceeds $K$ but does not include validation frames. In the held-out-object study, it also includes real data from the evaluation objects. On one RTX 6000 Ada GPU, aligner training takes about $20$ minutes, pseudo-labeling $0.5\,$s per frame, and student training $1.3$ hours for 50 epochs or $2.2$ hours for 150.

\noindent\textbf{Metrics.} We report all-pixel depth MSE (mm$^2$), contact MAE within the GT contact region (depth $>0$; mm), contact IoU (predicted depth $>0.05\,$mm versus GT depth $>0$), and mean rotation error (deg). Because roughly $90\%$ of the pixels are background, contact MAE and IoU expose errors that all-pixel MSE can obscure. Dense metrics use the checkpoint with the lowest validation depth MSE; sensor-A pose metrics use the checkpoint with the lowest validation rotation error.
